\subsection{Análisis de Requerimientos}

\subsubsection{Cuadro de relación de requerimientos y causas raíz}
En esta tabla se representará la relación de correspondencia entre los requerimientos funcionales y no funcionales con las causas raíz.

\begin{longtable}{p{3.5cm} p{4.5cm} p{4cm} p{2cm}}
    \caption{Relación entre los requerimientos y las causas raíz}
    \label{tab:trazabilidad_requerimientos} \\
    
    \toprule
    \textbf{Rama de Causa} &
     \textbf{Causa Raíz} &
      \textbf{Requerimiento(s) Asociado(s)} &
    \\ 
    \midrule
    \endfirsthead
    
    \multicolumn{4}{l}{\textit{(Continuación de la tabla anterior)}} \\
    \toprule
    \textbf{Rama de Causa} &
     \textbf{Causa Raíz} &
      \textbf{Requerimiento(s) Asociado(s)}\\ 
    \midrule
    \endhead
    
    \midrule
    \multicolumn{4}{r}{\textit{(Continúa en la siguiente página)}} \\
    \endfoot
    
    \bottomrule
    \endlastfoot

    Seguridad y trazabilidad de modificaciones &
     Permisos de consulta y modificación de datos laborales y salariales insuficientemente diferenciados según las responsabilidades del usuario. &
      M0-01, M0-03,RF-M1 \\
    \midrule
    Gestión de la información &
     Registros insuficientes de los cambios realizados en los datos laborales y resultado de la planilla. &
      M1-01, M1-02, M0-03, M1-04  \\
    &
     Información laboral y remunerativa almacenada en archivos separados entre RRHH y Contabilidad. &
      M4-01, M4-02, M1-03 \\
    &
     Circulación de distintas versiones de los archivos de planilla sin un control definido de actualización. &
      M5-04 \\
    \midrule
    Preparación de entrega de documentos &
     Transcripción manual de asistencias, permisos, horas extras y otras incidencias laborales de planillas. &
      M4-01, M4-02, M4-03, M5-04 \\
    \midrule
    Cálculo e incidencias &
     Modificación manual de fórmulas y parámetros de cálculo sin validación uniforme. &
      M4-05, M4-04, M3-01, M3-02, M3-03 \\
    &
     Preparación separada de información para pagos bancarios y declaraciones, mediante transferencias o registros manuales. &
      M1-01, M4-01, M4-03, M4-04 \\
    &
     Generación y distribución individual de boletas de pago mediante actividades manuales o semimanuales. &
      M5-01, M3-04 \\
\end{longtable}

\subsubsection{Requerimientos identificados del proyecto}
    \begin{itemize}
    \item RF-01: Aplicar acceso por Roles - El sistema debe autenticar a los usuarios mediante credenciales y asignar un token JWT, permitiendo al usuario ejecutar acciones segun su rol.
    \item RF-02: Registrar cambios en bicatora de auditoria - El sistema debe registrar automáticamente en la base de datos las acciones del usuario a nivel avanzado, permitiendo una reconstruccion de los datos.
    \item RF-03: Administrar legajos de los empleados - El sistema debe permitir gestionar la información maestra de los empleados.
    \item RF-04: Ingresar informacion de manera masiva - El sistema debe permitir la registrar informacion mediante archivos estructurados.
    \item RF-05: Gestionar contratos de los empleados - El sistema debe tenen la funcionabilidad de administrar los contratos laborales, notificar sobre el vencimientos, cese o fin de estos.
    \item RF-06: Gestionar las solicitudes del empleado - El sistema debe permitir hacer seguimiento a las solicitudes de los empleados, estos siendo aprobados o rechazados por el personal encargado de gestionar estos procedimientos. 
    \item RF-07: Asignar tareas a los usuarios - El sistema debe ofrecer al Contador asignar tareas operativas a los asistentes, monitoreando su estado y aprobar la tarea.
    \item RF-08: Cargar archivos estructurados y anotaciones - El sistema debe permitir al Asistente adjuntar documentos estructurados y comentarios en las tareas asignadas, así como solicitar prórrogas de entrega.
    \item RF-09: Modificar la variables ley - El sistema debe permitir al Contador actualizar el valor de las variables relacionadas para el calculo de planillas.
    \item RF-10: Distribución Masiva de Boletas por Correo - El sistema debe distribuir boletas electronicas a los correos del empleado una vez se haya aprobado el pago de planillas. .
    \item RF-11: Ingresar insidencias del mes - El sistema debe permitir cargar las incidencias del mes como vacaciones, horas extra, etc.
    \item RF-12: Generar Cálculo Masivo de Planillas  - El sistema debe ejecutar el cálculo masivo de planillas (general, filtrado o individual), aplicando descuentos, bonos, entre otros.
    \item RF-13: Cierre, Congelamiento Inmutable y Reapertura Excepcional - El sistema debe permitir al Contador ejecutar el "Cierre Definitivo" del periodo, congelando la base de datos para impedir mutaciones. Debe proveer un botón de "Reapertura Excepcional" que audite el desbloqueo en los Trail Logs.
    \item RF-14: Generar de Boletas electronicas - El sistema debe generar automáticamente las boletas de pago individuales en formato PDF tras el cierre de la planilla y guardarlas en el historial del legajo.
    \item RF-15: Reportes de Costos por Centro de Costo - El sistema debe exportar resúmenes financieros mensuales por áreas organizacionales.
\end{itemize}
\subsubsubsection{Requerimientos no Funcionales}